\documentclass[conference]{IEEEtran}
\IEEEoverridecommandlockouts

\usepackage{cite}
\usepackage{amsmath,amssymb,amsfonts}
\usepackage{graphicx}
\usepackage{textcomp}
\usepackage{xcolor}
\usepackage{booktabs}
\usepackage{url}
\usepackage{hyperref}

\hypersetup{
    colorlinks=true,
    linkcolor=black,
    citecolor=black,
    urlcolor=blue
}

\begin{document}

\title{TruthLens: Explainable Multi-Modal Detection of Synthetic Images, Audio and Text}

\author{\IEEEauthorblockN{Shashank Singh}
\IEEEauthorblockA{\url{https://github.com/Shashank17singh/TruthLens}}
}

\maketitle

\begin{abstract}
TruthLens is an open-source system that scores images, audio and text for signs of machine generation and explains each score. It pairs an RGB backbone with a 2D DCT frequency branch and Grad-CAM for images, a Light CNN over LFCC features for audio, and, for text, a logistic regression over GPT-2 statistics and a fine-tuned RoBERTa logit. Every dataset is split so that identities, speakers or prompts never cross between training and test. On the MAGE test set, the text ensemble reaches 0.997 AUC (95\% interval 0.997 to 0.998), detects 90.1\% of text from held-out LLaMA generators, and keeps 0.997 AUC on HC3 after removing rows that overlap MAGE. It still flags 61.8\% of TOEFL essays by non-native writers as AI-written, so it must not be used to judge individual writers. The LCNN audio model reaches 0.997 AUC on the ASVspoof 2019 LA validation set after three epochs on a laptop CPU; its evaluation-set results and the image detector's accuracy are not yet measured. All models export to ONNX, and the LCNN scores a 30-second clip in 280\,ms on four CPU threads.
\end{abstract}

\begin{IEEEkeywords}
Deepfake detection, machine-generated text, audio anti-spoofing, explainability, dataset leakage, CPU inference.
\end{IEEEkeywords}

\section{Introduction}

Image diffusion models, neural vocoders and large language models now produce media that people struggle to tell apart from real recordings and writing. Detectors for each modality exist, but three problems limit their use outside benchmarks. Most handle a single modality. Most return one score with no reason attached. Many are evaluated on splits where the same face, speaker or prompt appears in both training and test data, which inflates their results.

TruthLens addresses these problems in one codebase:
\begin{itemize}
    \item \textbf{Three detectors, one interface.} Images, audio and text each get a calibrated probability, a label and an explanation: a Grad-CAM map for images, a per-window timeline for audio, and per-sentence scores for text.
    \item \textbf{Leak-free splits.} A shared split engine keeps identities, speakers and prompt groups disjoint and holds out whole generators for an unseen-generator test.
    \item \textbf{Measured, reported limits.} The text detector's results include a bias test on essays by non-native writers and a robustness test against paraphrasing, mixing and truncation. Results that have not been measured are reported as missing.
    \item \textbf{CPU serving.} Models export to ONNX with output checks against PyTorch, and a FastAPI service returns results and signed PDF reports.
\end{itemize}

\section{Related Work}

\subsection{Visual Deepfake Detection}
Convolutional backbones such as ConvNeXt~\cite{liu2022convnet} and EfficientNet, and Vision Transformers, are the usual image detectors. Frank et al.~\cite{frank2020leveraging} showed that the upsampling in GANs leaves artifacts in the frequency spectrum, which motivates the DCT branch used here. FaceForensics++~\cite{rossler2019faceforensics} is the standard face-manipulation benchmark.

\subsection{Acoustic Anti-Spoofing}
ASVspoof 2019~\cite{todisco2019asvspoof} benchmarks synthetic and converted speech. LFCC front-ends with a Light CNN using Max-Feature-Map activations~\cite{wu2018lcnn} are a strong lightweight baseline. Self-supervised encoders such as wav2vec 2.0~\cite{baevski2020wav2vec} score higher but are large and slow on CPUs.

\subsection{Machine-Generated Text}
Zero-shot detectors use a language model's token probabilities: GLTR~\cite{gehrmann2019gltr} tracks the rank of each token, DetectGPT~\cite{mitchell2023detectgpt} probes the curvature of the log probability, and Binoculars~\cite{hans2024binoculars} contrasts two related models. Liang et al.~\cite{liang2023gpt} found that such detectors flag essays by non-native English writers far more often than essays by native writers.

\section{Datasets and Splits}

\begin{itemize}
    \item \textbf{ASVspoof 2019 LA}~\cite{todisco2019asvspoof}: 121,456 clips from 20 generators, with the official speaker-disjoint train, development and evaluation partitions.
    \item \textbf{MAGE}: human and machine text across ten domains and many generators. Splits are disjoint by prompt group, and the LLaMA generators are held out as an unseen test.
    \item \textbf{HC3}: human and ChatGPT answers. 497 rows open with the same words as MAGE training texts and are excluded from cross-dataset tests.
    \item \textbf{ChatGPT Detector Bias}~\cite{liang2023gpt}: 831 essays, including 89 TOEFL essays by non-native writers.
\end{itemize}

The split engine connects samples that share an identity, speaker or prompt into groups with a union-find pass, assigns whole groups to splits, and fails if any group crosses a split boundary.

\section{Method}

\subsection{Images: RGB and DCT Fusion}
An RGB backbone (EfficientNet-B0, ConvNeXt-Tiny or ViT-B/16) runs beside a frequency branch. The frequency branch converts the image to luminance and applies an orthonormal 2D type-II DCT inside the network:
\begin{equation}
\begin{split}
    D(u, v) = \alpha(u)\alpha(v) \sum_{i=0}^{N-1}\sum_{j=0}^{N-1} Y(i, j) &\cos\frac{\pi(2i+1)u}{2N} \\
    &\times \cos\frac{\pi(2j+1)v}{2N}
\end{split}
\end{equation}
A small CNN reads the log-magnitude spectrum, and a late-fusion layer combines both branches. Grad-CAM~\cite{selvaraju2017gradcam} on the last convolutional stage gives the explanation, checked with a model-randomization test~\cite{adebayo2018sanity}.

\subsection{Audio: LFCC and LCNN}
Audio is resampled to 16\,kHz mono and cut into 4-second windows with 50\% overlap. The front-end computes 20 LFCCs from 70 linear filters plus their deltas and delta-deltas, inside the network. A Light CNN with Max-Feature-Map activations classifies each window:
\begin{equation}
    \text{MFM}(x)_k = \max(x_k, x_{k + C/2}), \quad k \in [1, C/2]
\end{equation}
Window logits are reduced by both mean and max, each calibrated on validation data, which gives a per-window timeline for the interface.

\subsection{Text: Statistics and a Fine-Tuned Transformer}
GPT-2 (the performer) and DistilGPT-2 (the observer) score every token once, using 512-token windows. Eight features come from that pass: log perplexity, burstiness (the standard deviation of per-sentence log perplexity), the share of tokens ranked in the top 1, 10, 100 and 1000, mean entropy, and a Binoculars score~\cite{hans2024binoculars}. RoBERTa-base is fine-tuned on MAGE over 512-token windows, and its window-averaged logit is a ninth feature. A standardized logistic regression combines the nine features. Its threshold is the validation equal-error point. Sentence scores reuse the per-token values, and a bootstrap over sentences gives a 90\% band for each text.

\section{Implementation}

A FastAPI service queues jobs through Celery and Redis and stores media in MinIO, deleting it after 24 hours. Models are exported to ONNX and checked against PyTorch outputs. Finished analyses produce a PDF report whose digest is signed with HMAC-SHA256 and linked by QR code to a public verification page. A React web app and a Manifest V3 browser extension call the same API.

\section{Results}

\subsection{Detection}
The text models were trained on one Kaggle T4 GPU. RoBERTa-base was fine-tuned for four epochs of 3,000 steps, and its best validation AUC, 0.996, came at epoch 2. The logistic regression was trained on 4,999 MAGE texts. AUC intervals come from a bootstrap that resamples whole prompt groups. Table~\ref{tab:text} gives the results.

\begin{table}[t]
\centering
\caption{Text detection. AUC with 95\% grouped bootstrap interval.}
\label{tab:text}
\begin{tabular}{lrl}
\toprule
\textbf{Test set} & \textbf{Texts} & \textbf{AUC} \\
\midrule
\multicolumn{3}{l}{\textit{Statistics and RoBERTa}} \\
MAGE test & 3,999 & 0.997 (0.997--0.998) \\
MAGE, unseen LLaMA (AI only) & 1,999 & 90.1\% detected \\
HC3, overlap removed & 2,000 & 0.997 (0.995--0.998) \\
Detector bias set & 831 & 0.794 (0.761--0.825) \\
\midrule
\multicolumn{3}{l}{\textit{Statistics only}} \\
MAGE test & 3,999 & 0.782 (0.767--0.797) \\
\bottomrule
\end{tabular}
\end{table}

At its threshold, the ensemble reaches 97.6\% accuracy on the MAGE test set and 94.2\% on HC3, where it flags 9.3\% of human Reddit ELI5 answers. Adding the transformer raises test AUC from 0.782 to 0.997, so most of the detector's accuracy comes from RoBERTa rather than the statistics.

The LCNN with LFCC trained on the development laptop's CPU, capped at 300 steps per epoch. After three epochs its validation AUC was 0.997. The same validation set picked the checkpoint, so this figure is optimistic. Evaluation-set EER and min t-DCF have not yet been measured. The image detector has not been trained on real data, so no image accuracy is reported.

\subsection{Bias Against Non-Native Writers}
Table~\ref{tab:bias} gives the share of human essays that the ensemble flags as AI-written. It flags most TOEFL essays even though it rarely flags native student writing, which reproduces the bias Liang et al.~\cite{liang2023gpt} found in commercial detectors. The service shows a disclaimer that scores must not be used alone to decide academic misconduct.

\begin{table}[t]
\centering
\caption{Human essays flagged as AI-written}
\label{tab:bias}
\begin{tabular}{lrr}
\toprule
\textbf{Writers} & \textbf{Essays} & \textbf{Flagged} \\
\midrule
TOEFL, non-native writers & 89 & 61.8\% \\
College essays & 70 & 14.3\% \\
Stanford CS224N & 145 & 8.3\% \\
US 8th grade (Hewlett) & 88 & 5.7\% \\
\bottomrule
\end{tabular}
\end{table}

\subsection{Robustness of the Statistical Detector}
Table~\ref{tab:robust} tests the statistics-only model on 300 AI and 300 human MAGE test texts. Paraphrasing with a T5 model or replacing a quarter of the sentences with human writing each lowers detection from 71.3\% to about 47\%. The same test has not yet been run on the ensemble with RoBERTa.

\begin{table}[t]
\centering
\caption{Statistics-only detector under evasion}
\label{tab:robust}
\begin{tabular}{lr}
\toprule
\textbf{Condition} & \textbf{AI detected} \\
\midrule
Unmodified & 71.3\% \\
Paraphrased by T5 & 47.7\% \\
25\% of sentences human & 47.0\% \\
50\% of sentences human & 37.3\% \\
First 75 words & 53.7\% \\
First 50 words & 37.3\% \\
\midrule
Human texts flagged (unmodified) & 30.7\% \\
\bottomrule
\end{tabular}
\end{table}

\subsection{CPU Latency}
Table~\ref{tab:latency} gives median ONNX Runtime latency on an Intel Core i5-10210U laptop. Image models score one 224$\times$224 image; audio models score a 30-second clip. Latency depends on the architecture, not on the trained weights, so these figures hold for the trained models.

\begin{table}[t]
\centering
\caption{Median ONNX latency on an Intel Core i5-10210U}
\label{tab:latency}
\begin{tabular}{lrrr}
\toprule
\textbf{Model} & \textbf{MB} & \textbf{1 thread} & \textbf{4 threads} \\
\midrule
EfficientNet-B0 + DCT & 18.4 & 45\,ms & 28\,ms \\
ConvNeXt-Tiny + DCT & 113.7 & 177\,ms & 93\,ms \\
ViT-B/16 + DCT & 345.9 & 553\,ms & 280\,ms \\
\midrule
LCNN, LFCC & 2.2 & 533\,ms & 280\,ms \\
ResNet-18, LFCC & 45.9 & 940\,ms & 640\,ms \\
wav2vec2 & 378.9 & 30.5\,s & 9.8\,s \\
\bottomrule
\end{tabular}
\end{table}

\section{Limitations}

\begin{enumerate}
    \item \textbf{Missing results.} The image detector has no accuracy result on real data, and the audio detector has only a validation AUC. The DeBERTa-v3-small run failed on its first attempt and is being repeated.
    \item \textbf{Bias.} The text detector flags most essays by non-native writers. It must not be used to judge individual writers.
    \item \textbf{Evasion.} Paraphrasing and mixing in human sentences defeat the statistical features. The ensemble's robustness has not been measured.
    \item \textbf{Domain.} MAGE and HC3 share sources and generators, so the high cross-dataset AUC may not hold for newer models or other kinds of writing.
\end{enumerate}

\section{Conclusion}

TruthLens provides explainable image, audio and text detection with leak-free splits and CPU serving. Fine-tuning RoBERTa lifts the text detector from 0.782 to 0.997 AUC on MAGE and keeps it at 0.997 on HC3, but its bias against non-native writers remains. The next steps are training the image detector on real data, evaluating the audio detector on the ASVspoof evaluation set with min t-DCF, and measuring the ensemble's robustness to paraphrasing.

\bibliographystyle{IEEEtran}
\bibliography{references}

\end{document}
